\section{Background}\label{sec:background}

\subsection{Random parameters in crash frequency analysis}

Crash counts are overdispersed relative to the Poisson distribution, and the negative binomial (NB) model, which mixes the Poisson rate over a gamma distribution, has long been the default for safety performance functions \citep{lord2010statistical}. Fixed parameters NB models assume that a covariate has the same effect at every site. Random parameters models relax this by letting coefficients vary across observations, which captures heterogeneity due to unobserved driver, vehicle, roadway and environmental factors \citep{mannering2016unobserved}. Following \citet{anastasopoulos2009note} and \citet{milton2008highway}, random parameters count and discrete outcome models have been applied to segment and intersection crash frequencies, crash rates and injury severities, often with extensions to heterogeneity in means and variances, correlated random parameters and temporal instability \citep{mannering2014analytic, mannering2018temporal}. Estimation is typically by simulated maximum likelihood with Halton draws \citep{bhat2003simulation, train2009discrete} or by full Bayesian methods \citep{elbasyouny2009accident}. \citet{buddhavarapu2016modeling} used P\'olya-Gamma augmentation to estimate finite mixture random parameters models for spatially correlated crash counts, which shows that the computational route used in this paper is practical for safety data.

\subsection{How specifications are chosen}

Textbook treatments describe a sequential procedure \citep{washington2020statistical}. Candidate covariates are first screened in a fixed parameters model. Each retained coefficient is then specified as random, and the random specification is kept when the standard deviation of the parameter distribution is statistically different from zero, otherwise the parameter is treated as fixed. In Bayesian applications the same decision is made by inspecting a credible interval for the standard deviation, or by comparing deviance information criteria of models with and without the random parameter \citep[e.g.,][]{zeng2017multivariate}. Because the number of candidate specifications grows as $3^p$, and faster once heterogeneity in means or correlation among random parameters is allowed, \citet{ahern2024extensive} formulated specification search as a combinatorial optimization problem and solved it with metaheuristics that minimize the Bayesian information criterion.

These procedures share two features. They return one specification, so the uncertainty about which parameters are random is not propagated into estimates of marginal effects or into decisions based on them. They also perform many tests of variance components, a setting in which classical tests are non-standard because the null value $\sigma_j=0$ lies on the boundary of the parameter space, and in which no adjustment is made for the number of tests.

\subsection{Bayesian variable selection and selection of random effects}

Spike-and-slab priors place a mixture of a point mass (or a narrow distribution) at zero and a diffuse distribution on each regression coefficient, so that the posterior probability of the slab component measures the evidence that a variable belongs in the model \citep{mitchell1988bayesian, george1993variable, ishwaran2005spike}. Treating the prior inclusion probability as unknown, with a beta hyperprior, yields an automatic multiplicity adjustment \citep{scott2010bayes}, and the median probability model is optimal for prediction under conditions given by \citet{barbieri2004optimal}. Continuous shrinkage alternatives include the Bayesian lasso \citep{park2008bayesian, kyung2010penalized}, the horseshoe \citep{carvalho2010horseshoe} and the spike-and-slab lasso \citep{rockova2018spike}. \citet{ohara2009review} compare several of these methods in practice.

Selecting random effects is harder than selecting fixed effects because a variance cannot be negative, so a spike at zero sits on the boundary of the parameter space. \citet{chen2003random} and \citet{kinney2007fixed} selected random effects in linear and logistic mixed models by reparameterizing the covariance matrix of the random effects. \citet{fruhwirth2008bayesian} and \citet{fruhwirth2010stochastic} introduced the device that we adopt: write the random component as $\sigma_j z_{ij}$ with $\sigma_j\in\mathbb{R}$, so that $\sigma_j$ enters the linear predictor linearly given $z_{ij}$ and can receive an ordinary spike-and-slab prior. A Gaussian slab on the signed $\sigma_j$ implies a scaled $\chi^2_1$ prior on the variance $\sigma_j^2$, which has substantial mass near zero. Inverse-gamma priors on variances, by contrast, have vanishing density at zero and can overstate heterogeneity when the data are weakly informative \citep{gelman2006prior}. \citet{bitto2019achieving} extended the idea to continuous shrinkage in time-varying parameter models.

Deciding jointly whether a coefficient enters and whether it varies is a bi-level selection problem in which the mean and the standard deviation of a parameter form a natural group. Bi-level and group spike-and-slab priors have been studied for linear regression \citep{xu2015bayesian, yang2020consistent, bai2022spike}, and heredity priors that let the inclusion of one term depend on another go back to \citet{chipman1996bayesian}. To our knowledge these tools have not been brought together for the specification of random parameters crash models, where the groups have a specific structure (a mean and a standard deviation), the likelihood is negative binomial, and the random components are observation specific rather than shared within clusters. \tbd{Confirm with a final literature check before submission, including recent AMAR, AA\&P and TR-B issues.}
